% !TeX root = ../main.tex

\setproblem{F}{\ProblemTitleF}
\section{\texorpdfstring{\pletter. \ptitle}{문제 \pletter}}

\begin{frame}[label=problem-F]
    \sectiontitle{\pletter}{\ptitle}
    \sectionmeta{
        출제자: \idtext{areian13}\\
        태그: \tagtext{\#implementation, \#simulation}\\
        예상 난이도: \tiertext{Bronze2}
    }
    \begin{itemize}
        \item 제출 17번, 정답 10명 (정답률 58.8\%)
        \item 처음 푼 사람: ㅅㅇㅅ (00:13:26)
    \end{itemize}
\end{frame}

\begin{frame}{\textbf{\pletter}. \ptitle}
    \textbf{동시에 충전하는 노트북 수 세기}
\begin{itemize}
        \item 어떤 시각에 노트북 $x$대가 동시에 충전 중이면 콘센트도 적어도 $x$개 필요합니다.
        \item 따라서 시각별로 충전 중인 노트북 수를 세어 최댓값을 구해 봅시다.
        \item 시각은 $0$부터 $10\,000$까지이므로 배열로 직접 셀 수 있습니다.
        \item 충전 일정이 $(s_i,e_i)$라면 $s_i,s_i+1,\ldots,e_i-1$번 칸을 각각 $1$씩 늘립니다.
        \item $e_i$에는 이미 충전이 끝났으므로 표시하지 않습니다.
        \item $(1,3)$과 $(3,5)$는 동시에 충전하지 않으므로 콘센트 하나를 함께 쓸 수 있습니다.
    \end{itemize}
\end{frame}

\begin{frame}{\textbf{\pletter}. \ptitle}
    \textbf{왜 최댓값만큼이면 충분할까?}
\begin{itemize}
        \item 동시에 충전하는 노트북 수의 최댓값을 $p$라 합시다. 콘센트가 $p$개보다 적으면 부족합니다.
        \item 반대로 $p$개가 있다면, 충전을 시작하는 노트북에 비어 있는 콘센트를 배정하면 됩니다.
        \item 빈 콘센트가 없다면 새 노트북까지 $p+1$대가 동시에 충전한다는 뜻이므로 모순입니다.
        \item 따라서 구한 최댓값 $p$가 필요한 콘센트의 최소 개수입니다.
        \item 가능한 시각의 수를 $T$라 하면 시간 복잡도는 $O(NT)$, 공간 복잡도는 $O(T)$입니다.
        \item 여기서는 $NT\le10^7$이므로 모든 일정의 구간을 직접 표시해도 됩니다.
    \end{itemize}
\end{frame}
